% ------------------------------------------------------------
\subsection{Cobertura de los escenarios}

Para comprobar que los escenarios no dejan fuera nada importante del análisis de contexto, se revisa qué escenarios cubren cada stakeholder, concern, restricción y tensión.

\textbf{Stakeholders.} Los 19 stakeholders aparecen en al menos un escenario.

\begin{tabla}{|L{1.5cm}|Y|L{1.5cm}|Y|}
\hline
\filacab \cab{STK} & \cab{Escenarios} & \cab{STK} & \cab{Escenarios} \\ \hline
\endhead
STK-01 & ESC-01, 02, 03, 04, 07, 08, 09, 10, 16 & STK-11 & ESC-11 \\ \hline
STK-02 & ESC-03, 04, 06 & STK-12 & ESC-07 (si Q-06 incluye revisión humana) \\ \hline
STK-03 & ESC-13, 17 & STK-13 & ESC-09, 10 \\ \hline
STK-04 & ESC-13, 14, 15 & STK-14 & ESC-08, 09, 10, 14, 18 \\ \hline
STK-05 & ESC-05, 12, 14, 16 & STK-15 & ESC-03, 05, 06, 07, 16 \\ \hline
STK-06 & ESC-12 & STK-16 & ESC-17 \\ \hline
STK-07 & ESC-13 & STK-17 & ESC-04, 07, 15 \\ \hline
STK-08 & ESC-11, 15 & STK-18 & ESC-15 \\ \hline
STK-09 & ESC-01, 02, 03, 11 & STK-19 & ESC-04, 18 \\ \hline
STK-10 & ESC-05, 12, 17 & & \\ \hline
\end{tabla}

\textbf{Concerns.} 25 de los 26 concerns tienen al menos un escenario. CRN-09 queda sin escenario de forma intencional, por el motivo que se explica al final de esta sección.

\begin{tabla}{|L{1.5cm}|Y|L{1.5cm}|Y|}
\hline
\filacab \cab{CRN} & \cab{Escenarios} & \cab{CRN} & \cab{Escenarios} \\ \hline
\endhead
CRN-01 & ESC-02, 03 & CRN-14 & ESC-13, 14, 15 \\ \hline
CRN-02 & ESC-01 & CRN-15 & ESC-11 \\ \hline
CRN-03 & ESC-12 (parcial) & CRN-16 & ESC-11 \\ \hline
CRN-04 & ESC-04 & CRN-17 & ESC-17 \\ \hline
CRN-05 & ESC-03, 06, 07 & CRN-18 & ESC-07 \\ \hline
CRN-06 & ESC-08 & CRN-19 & ESC-09 \\ \hline
CRN-07 & ESC-13 & CRN-20 & ESC-14 \\ \hline
CRN-08 & ESC-17 & CRN-21 & ESC-10 \\ \hline
CRN-09 & Sin escenario (documentación) & CRN-22 & ESC-03, 06 \\ \hline
CRN-10 & ESC-05 & CRN-23 & ESC-05 \\ \hline
CRN-11 & ESC-05 & CRN-24 & ESC-15 \\ \hline
CRN-12 & ESC-12, 17 & CRN-25 & ESC-04, 18 \\ \hline
CRN-13 & ESC-12 & CRN-26 & ESC-16 \\ \hline
\end{tabla}

\textbf{Restricciones y tensiones.}

\begin{tabla}{|L{1.5cm}|Y|L{1.5cm}|Y|}
\hline
\filacab \cab{ID} & \cab{Escenarios} & \cab{ID} & \cab{Escenarios} \\ \hline
\endhead
CON-01 & Sin escenario propio (ver CON-16) & CON-12 & ESC-02, 03 \\ \hline
CON-02 & ESC-01, 06, 08, 14 & CON-13 & ESC-13, 17 \\ \hline
CON-03 & ESC-05, 06 & CON-14 & ESC-13, 17 \\ \hline
CON-04 & ESC-09, 11 & CON-15 & ESC-12, 15 \\ \hline
CON-05 & ESC-11 & CON-16 & ESC-18 \\ \hline
CON-06 & ESC-01, 05, 08 & CON-17 & ESC-08, 14 \\ \hline
CON-07 & ESC-01 & CON-18 & ESC-15 \\ \hline
CON-08 & ESC-03, 06, 16 & TEN-01 & ESC-01, 05, 06, 12 \\ \hline
CON-09 & ESC-04 & TEN-02 & ESC-04, 07 \\ \hline
CON-10 & ESC-07 & TEN-03 & ESC-03, 16 \\ \hline
CON-11 & ESC-02, 04, 11 & TEN-04 & ESC-13, 14 \\ \hline
 & & TEN-05 & ESC-02, 11 \\ \hline
 & & TEN-06 & ESC-01, 09 \\ \hline
\end{tabla}

\textbf{Elementos sin escenario propio.}

\begin{itemize}
  \item \textbf{CRN-09} (justificar cada decisión estructural). Es una necesidad del proceso y de la documentación, no una propiedad del juego en ejecución ni de su código. Se atiende manteniendo la matriz de trazabilidad y registrando cada decisión con su motivo. Forzar un escenario de producto para este concern produciría una medida artificial.
  \item \textbf{CON-01} (desarrollar en C\#). Es una decisión de tecnología, no una situación que el sistema enfrente. Su efecto sobre la calidad llega a través de su restricción derivada, CON-16, que sí tiene escenario (ESC-18).
  \item \textbf{CRN-03} queda cubierto de forma parcial. ESC-12 cubre la parte de mantener el interés con variantes de reglas. La parte de aprender en pocos minutos se relaciona con ESC-02, pero no tiene una medida propia de tiempo de aprendizaje. Si el equipo de experiencia de usuario lo considera necesario, es el primer escenario que convendría agregar.
\end{itemize}

\textbf{Cómo se usarán los escenarios.} Los escenarios de prioridad (A, A), ESC-01, ESC-03, ESC-05, ESC-06, ESC-08 y ESC-09, son los que deben guiar la evaluación de las alternativas de arquitectura: una alternativa que no pueda cumplirlos se descarta aunque sea mejor en otros aspectos. Los de crecimiento sirven para decidir qué abstraer y qué no, de acuerdo con TEN-04. Cada medida se convierte, al construir el sistema, en una prueba del incremento semanal correspondiente.
